\documentclass[11pt,a4paper]{article}
\usepackage[utf8]{inputenc}
\usepackage[T1]{fontenc}
\usepackage[spanish]{babel}
\usepackage{hyperref}
\usepackage{geometry}
\geometry{margin=2.5cm}
\usepackage{parskip}

\title{El nombre}
\author{Alejandro Valenzuela Rivera\\\small Universidad Autónoma de Nuevo León}
\date{\today}
\begin{document}
\maketitle

\begin{quote}
\textbf{Sesgo declarado.} Este trabajo se sostiene desde el sesgo declarado de su autor, y no finge neutralidad. Lo que afirma se puede tumbar, y ese es el punto: el falsificador está escrito, no escondido.
\end{quote}

\section{La tesis}
\textbf{\textquestiondown{}Cómo se escribe una idea que no es de nadie todavía?}

\section{Por qué se sostiene}
\textbf{Porque su contexto refutable está escrito —`FUENTES`, `ATTRIBUTION`— y se puede tumbar sin que esta pregunta caiga.} \textbf{Y porque cada nombre del repo tiene su acto y su origen, y se puede comprobar leyendo.}

\section{De dónde viene}
\textbf{De `Tesis/ILUMINAR.md`, de `FUENTES.md` y `ATTRIBUTION.md`.} \textbf{Y de la autoridad: «si te sabes la fuente métela o triangula si es posible, porque esto me puede meter en problemas si no los citas bien».}

\section{Lo que queda abierto}
\textbf{\textquestiondown{}Cuándo un nombre prestado se vuelve un problema, y cuándo solo es un préstamo?}

\section{El caso}
\textbf{La Source Wall, de Jack Kirby (DC Comics, 1971): el borde que lo firma alguien que no es el que lo cruza.} \textbf{Y `AM`, de Harlan Ellison (1967): el nombre que se volvió ley sin autor.}

\section{Método}
\textbf{\textquestiondown{}Cada nombre del repo dice un acto? \textquestiondown{}Cada obra citada está en la tabla con su autoría y dónde se usa? \textquestiondown{}Algún nombre de obra se usa como nombre propio de un sistema? \textquestiondown{}Se puede citar sin preguntarle al autor?}

\section{Límites, y qué la tumba}
\textbf{Un nombre sin acto; una obra citada sin su referencia; un nombre de obra usado como nombre propio; una idea presentada como propia sin declarar de dónde viene.}

\section{Estado de la tesis}
\textbf{escrita.} \textbf{Falta el caso de alguien que reclame un nombre y se resuelva con `FUENTES` y no con un abogado.} ---

\section*{Cómo se leyó esta tesis}
Esto se escribió en el repositorio man-made-empathy, y la tesis está completa en vision del autor/tesis/man made empathy/el nombre/. La máquina que la compuso no agregó nada que no estuviera escrito: si algo de esto es falso, se dice en el mismo sitio, y se corrige con fecha.
\end{document}
